% Handschrift: Seiten 8--13; Betragsregeln von Seite 6 hier eingeordnet.
\section{Rechnen in kartesischer Darstellung}
\label{sec:complex_numbers_cartesian_arithmetic}
Wir rechnen mit den vertrauten algebraischen Gesetzen und ersetzen dabei
$i^2$ durch $-1$. Addition und Subtraktion lassen sich unmittelbar in der
Zahlenebene deuten; die Geometrie der Multiplikation folgt später.

\begin{KeyPointBox}{Addition und Subtraktion}{complex_numbers_addition}
  Für $z_1=x_1+iy_1$ und $z_2=x_2+iy_2$ werden die beiden Anteile
  getrennt addiert beziehungsweise subtrahiert:
  \[
    z_1\pm z_2=(x_1\pm x_2)+i(y_1\pm y_2).
  \]
  Geometrisch entspricht dies der Addition beziehungsweise Subtraktion
  der zugehörigen Pfeile.
\end{KeyPointBox}
\begin{ExampleBox}{Summe und Differenz in der Zahlenebene}{complex_numbers_sum_difference}
  Bestimmen Sie für $z_1=2+i$ und $z_2=3-2i$ die Summe und die Differenz.
  \[
    \begin{aligned}
      z_1+z_2&=(2+3)+i(1-2)=\Cloze{3cm}{5-i},\\[5pt]
      z_1-z_2&=(2-3)+i(1+2)=\Cloze{3cm}{-1+3i}.
    \end{aligned}
  \]
  \begin{center}
    \CachedTikzFigure{content/resources/chapters/complex_numbers/tikz/complex_addition_subtraction.tex}
  \end{center}
\end{ExampleBox}
\clearpage
Für eine Zahl und ihre konjugierte Zahl heben sich beim Addieren die
Imaginärteile und beim Subtrahieren die Realteile auf:
\[
  \begin{aligned}
    z+z^*&=(x+iy)+(x-iy)=2x
      &&\Longrightarrow& \operatorname{Re}(z)&=\frac{z+z^*}{2},\\[6pt]
    z-z^*&=(x+iy)-(x-iy)=2iy
      &&\Longrightarrow& \operatorname{Im}(z)&=\frac{z-z^*}{2i}.
  \end{aligned}
\]
\begin{ExampleBox}{Komponenten durch Konjugation zurückgewinnen}{complex_numbers_components_conjugation}
  Prüfen Sie beide Formeln an $z=3+2i$ und $z^*=3-2i$.
  \[
    \operatorname{Re}(z)=\frac{(3+2i)+(3-2i)}2=\Cloze{1.5cm}{3},\qquad
    \operatorname{Im}(z)=\frac{(3+2i)-(3-2i)}{2i}=\Cloze{1.5cm}{2}.
  \]
\end{ExampleBox}

\subsection{Multiplikation}
Ausmultiplizieren und Zusammenfassen nach Real- und Imaginärteil liefert
\[
  \begin{aligned}
    z_1z_2&=(x_1+iy_1)(x_2+iy_2)\\[5pt]
      &=x_1x_2+ix_1y_2+iy_1x_2+i^2y_1y_2\\[5pt]
      &=(x_1x_2-y_1y_2)+i(x_1y_2+y_1x_2).
  \end{aligned}
\]
\begin{RemarkBox}
  Real- und Imaginärteil eines Produkts sind im Allgemeinen \emph{nicht}
  die Produkte der einzelnen Real- und Imaginärteile. Die gemischten
  Terme und $i^2=-1$ müssen berücksichtigt werden.
\end{RemarkBox}
\begin{KeyPointBox}{Produkt mit der konjugierten Zahl}{complex_numbers_conjugate_product}
  Beim Multiplizieren einer Zahl mit ihrer konjugierten Zahl verschwinden
  die gemischten Terme:
  \[
    zz^*=(x+iy)(x-iy)=x^2-i^2y^2=x^2+y^2=|z|^2\in\R_{\ge0}.
  \]
  Damit lässt sich der Betrag auch durch $|z|=\sqrt{zz^*}$ ausdrücken.
\end{KeyPointBox}
\begin{ExampleBox}{Den Betrag durch ein Produkt prüfen}{complex_numbers_modulus_product_example}
  Berechnen Sie für $z=3+2i$ das Produkt $zz^*$ und vergleichen Sie
  mit der geometrisch bestimmten Länge $\sqrt{13}$.
  \begin{ClozeArea}{4}
    \[
      (3+2i)(3-2i)=9-(2i)^2=9+4=13,
      \qquad |z|=\sqrt{13}.
    \]
  \end{ClozeArea}
\end{ExampleBox}

\clearpage
\subsection{Division}
Um einen Quotienten kartesisch darzustellen, machen wir seinen Nenner
mithilfe der komplexen Konjugation reell.
\begin{KeyPointBox}{Erweitern mit dem konjugierten Nenner}{complex_numbers_division}
  Für $z_2\ne0$ gilt mit $z_2z_2^*=|z_2|^2$ die Umformung
  \[
    \begin{aligned}
      \frac{z_1}{z_2}
        &=\frac{z_1z_2^*}{|z_2|^2}
         =\frac{(x_1+iy_1)(x_2-iy_2)}{x_2^2+y_2^2}\\[7pt]
        &=\frac{x_1x_2+y_1y_2}{x_2^2+y_2^2}
          +i\,\frac{y_1x_2-x_1y_2}{x_2^2+y_2^2}.
    \end{aligned}
  \]
  Auch beim Dividieren dürfen Real- und Imaginärteile nicht getrennt
  durcheinander geteilt werden.
\end{KeyPointBox}

\begin{ExampleBox}{Einen Quotienten kartesisch darstellen}{complex_numbers_quotient_example}
  Bringen Sie $(3+2i)/(1-i)$ auf die Form $x+iy$, indem Sie den Nenner reell machen.
  \begin{ClozeArea}{6}
    \[
      \begin{aligned}
        \frac{3+2i}{1-i}
          &=\frac{(3+2i)(1+i)}{(1-i)(1+i)}\\[6pt]
          &=\frac{3+5i+2i^2}{1-i^2}
           =\frac{1+5i}{2}=\frac12+\frac52i.
      \end{aligned}
    \]
  \end{ClozeArea}
  Realteil und Imaginärteil des Quotienten sind also $1/2$ und $5/2$.
\end{ExampleBox}
\begin{KeyPointBox}{Die imaginäre Einheit im Nenner}{complex_numbers_reciprocal_i}
  Beim Wechsel von $i$ aus dem Nenner in den Zähler entsteht ein Minuszeichen:
  \[
    \frac1i=\frac{-i}{i(-i)}=\frac{-i}{1}=-i,
    \qquad \frac{a}{i}=-ia\quad(a\in\C).
  \]
\end{KeyPointBox}

\clearpage
\subsection{Rechengesetze}
Die aus dem Reellen bekannten Gesetze gelten auch für komplexe Zahlen:
\[
  \begin{array}{@{}ll@{}}
    z_1+z_2=z_2+z_1,\quad z_1z_2=z_2z_1
      &\text{Kommutativgesetze},\\[7pt]
    z_1+(z_2+z_3)=(z_1+z_2)+z_3
      &\text{Assoziativgesetze},\\[7pt]
    z_1(z_2z_3)=(z_1z_2)z_3&\\[7pt]
    z_1(z_2+z_3)=z_1z_2+z_1z_3
      &\text{Distributivgesetz}.
  \end{array}
\]
Für die Konjugation ergeben sich außerdem die Beziehungen
\[
  (z_1+z_2)^*=z_1^*+z_2^*,\quad
  (z_1z_2)^*=z_1^*z_2^*,\quad
  \left(\frac{z_1}{z_2}\right)^*=\frac{z_1^*}{z_2^*}\quad(z_2\ne0).
\]
\begin{ExampleBox}{Die Reihenfolge des Zusammenfassens wählen}{complex_numbers_associativity}
  Berechnen Sie $(1+i)(3+2i)(3-2i)$ mit zwei verschiedenen Klammerungen.
  Nutzen Sie im ersten Weg das konjugierte Paar.
  \begin{ClozeArea}{7}
    \[
      \begin{aligned}
        (1+i)\bigl((3+2i)(3-2i)\bigr)
          &=(1+i)\cdot13=13+13i,\\[6pt]
        \bigl((1+i)(3+2i)\bigr)(3-2i)
          &=(1+5i)(3-2i)\\[6pt]
          &=3-2i+15i-10i^2=13+13i.
      \end{aligned}
    \]
  \end{ClozeArea}
\end{ExampleBox}

\begin{KeyPointBox}{Rechenregeln für den Betrag}{complex_numbers_modulus_laws}
  Für komplexe Zahlen gelten die Produkt- und Quotientenregel sowie
  die \emph{Dreiecksungleichung} \EnglishTerm{triangle inequality}:
  \[
    |z|\ge0,\qquad |z_1z_2|=|z_1||z_2|,\qquad
    \left|\frac{z_1}{z_2}\right|=\frac{|z_1|}{|z_2|}\quad(z_2\ne0),
  \]
  \[
    |z_1+z_2|\le|z_1|+|z_2|.
  \]
\end{KeyPointBox}
